\section*{Chapter \ref{ch:s1:alternative-data-strategies} --- Alternative Data Strategies}

\begin{solution}{exo:s1:alternative-data-strategies:1}
With noise twice the surprise's spread the correlation is $1/\sqrt{1 + 2^2} = 0.45$, and the best slope is $1/(1 + 2^2) = 0.2$: the \omterm{def:s1:alternative-data-strategies:nowcast}{nowcast} shrinks the reading
by four fifths. The chapter's fitted values are 0.46 and 0.21.
\end{solution}

\begin{solution}{exo:s1:alternative-data-strategies:2}
The \omterm{def:s1:alternative-data-strategies:nowcast}{nowcast} is $0.2 \times 2 = 0.4$ standard deviations of surprise, and the expected reaction $3 \times 1.6\% \times 0.4 = 1.92\%$. Half the market moves
$0.5 \times 1.92\% = 0.96\%$ of it on delivery, leaving the other half for the report.
\end{solution}

\begin{solution}{exo:s1:alternative-data-strategies:3}
Post-publication decay follows the spread of an idea that anyone can copy at no cost; \omterm{def:s1:alternative-data-strategies:decay}{data decay} follows the number of buyers of a dataset that costs money
and may be sold exclusively. Its pace can be slowed by exclusivity and speed, and it can be estimated from the dataset's sales rather than from the calendar.
\end{solution}

\begin{solution}{exo:s1:alternative-data-strategies:4}
The book turns over 36 times its gross a year, one way, at ten basis points per unit traded; when most of the reaction happens on delivery, the gross return falls below the
costs before the share with the data reaches one: at 0.9 the gross Sharpe ratio is 0.15 and the net $-0.43$.
\end{solution}

\begin{solution}{exo:s1:alternative-data-strategies:5}
With the data a day earlier the firm is positioned before the others trade on it, so it earns their move on delivery as well as the rest at the report:
the reaction's total is unchanged, and the firm holds through all of it, whatever the share of the market that follows.
\end{solution}

\begin{solution}{exo:s1:alternative-data-strategies:6}
Measure the price move on delivery days (or in the days before reports) per unit of the \omterm{def:s1:alternative-data-strategies:nowcast}{nowcast}: with no one else using the data it is zero; as usage
grows it rises toward the whole expected reaction. The ratio of the pre-report move to the total estimates the share.
\end{solution}

\begin{solution}{exo:s1:alternative-data-strategies:7}
With noise three times the spread (a correlation of 0.34) the book at half diffusion earns a Sharpe ratio of 1.35 and 17.6\% a year, against 1.86 and 24.2\% with
noise twice the spread; without diffusion 3.20 against 4.41. A noisier panel predicts less of each surprise, so every position carries less of the reaction
and the same costs.
\end{solution}

\begin{solution}{exo:s1:alternative-data-strategies:8}
A vendor's history may be backfilled (reconstructed after the fact and fitted to the outcomes), and it does not show how many other clients now have the
data. Demand point-in-time deliveries, trial the data live for a period, and estimate its current diffusion from the pre-report moves.
\end{solution}

\begin{solution}{pb:s1:alternative-data-strategies:1}
\begin{enumerate}
\item An estimate of a later-reported quantity from earlier data; a strategy trading on such \omterm{def:s1:alternative-data-strategies:nowcast}{nowcasts} from non-traditional data.
\item Map the data to companies, \omterm{def:s1:alternative-data-strategies:nowcast}{nowcast} the reported quantity, translate it into an expected reaction, and hold over the window in which the market learns.
\item A tenth of the companies, readings with noise twice the surprise's spread (correlation 0.46), delivered ten days early; slope 0.21 from the first two
years.
\item Nowcast-proportional positions from the delivery's close to the announcement's close, gross one.
\item Satellite data let sophisticated investors target bad reports, with more informed shorting and lower liquidity: more asymmetry, not immediately better
prices.
\item Improving employer ratings outperform declining ones and forecast earnings surprises.
\item Prices ignore news about principal customers; a supplier strategy earned over 150 basis points a month.
\item Alternative data raised price informativeness by lowering information costs.
\item The loss of value as more investors trade the same data earlier; a share $\phi$ moves $\phi$ of the expected reaction on delivery.
\item 4.41, 3.20, 1.86, 0.45 and $-0.43$ for shares of 0, 0.25, 0.5, 0.75 and 0.9.
\item Each buyer moves an equal part of the reaction to the delivery day, out of the book's reach.
\item It keeps the Sharpe ratio near 4.3--4.4 at every level of diffusion.
\item \textbf{Named result.} The \omterm{def:s1:alternative-data-strategies:nowcast}{nowcast} book's Sharpe ratio after costs falls from 4.41 when no one else has the panel to 1.86 when half the market has it and
$-0.43$ at nine tenths, while a firm that receives the panel a day earlier keeps 4.28 to 4.43.
\item Track the pre-report move per unit of \omterm{def:s1:alternative-data-strategies:nowcast}{nowcast} and the book's realised reaction per trade.
\item Late or changed deliveries, coverage losses, revisions, vendor failures and licence terms.
\item Insist on point-in-time history, check for backfill, and trial it live.
\item What it adds to the firm's returns after costs over its expected life, given its diffusion.
\item The revenue \omterm{def:s1:alternative-data-strategies:nowcast}{nowcast} trade, whose data are sold to many buyers.
\item They trade the same reports before they happen; chapter~7 trades them after.
\item Time: knowledge of a quantity before the market has it.
\end{enumerate}
\end{solution}

\begin{solution}{iq:s1:alternative-data-strategies:1}
Coverage and history (point in time?), the mapping to companies, the correlation with the reported quantity on data the vendor could not have fitted, the
incremental IC over the firm's existing signals, the lag, the cost, who else buys it, and the legal terms.
\end{solution}

\begin{solution}{iq:s1:alternative-data-strategies:2}
History reconstructed after the fact, often fitted to known outcomes, so the backtest looks better than live use. Detect it by comparing fit before and
after the vendor's launch date, by coverage jumps and level shifts at the launch, and by asking for delivery timestamps (Book~7, chapter~12).
\end{solution}

\begin{solution}{iq:s1:alternative-data-strategies:3}
More buyers (diffusion), a change in the panel (provider, coverage, methodology), a change in what the companies report, or crowding by similar datasets.
Check the pre-report move per unit of signal and the panel's own statistics.
\end{solution}

\begin{solution}{iq:s1:alternative-data-strategies:4}
Ingest and validate the delivery (schema, coverage, ranges against history), map to company identifiers point in time, compute \omterm{def:s1:alternative-data-strategies:nowcast}{nowcasts} with the current
model, pass them to the optimiser with the other signals, generate orders for the close, and log everything for reproduction.
\end{solution}

\begin{solution}{iq:s1:alternative-data-strategies:5}
Material non-public information (data obtained in breach of a duty), personal data and privacy law, licence terms and exclusivity, and web-scraping terms of
service; every dataset needs a compliance review of how it was collected.
\end{solution}

\begin{solution}{iq:s1:alternative-data-strategies:6}
With expected reaction $J$ per unit of \omterm{def:s1:alternative-data-strategies:nowcast}{nowcast}, a trader at delivery earns $(1 - \phi)J$ per unit, since $\phi J$ moves when the others trade. A trader with the
data $k$ days earlier, holding through, earns the whole $J$ (the others' move included), less any drift of the extra days; the value of speed is $\phi J$.
\end{solution}
